% BHU PYQ -- Auto-generated & error-corrected
\documentclass[11pt,a4paper]{article}

\usepackage[a4paper,top=2.5cm,bottom=2.5cm,left=2.8cm,right=2.8cm,headheight=14pt]{geometry}
\usepackage{amsmath,amssymb}
\usepackage[T1]{fontenc}
\usepackage[utf8]{inputenc}
\usepackage{lmodern}
\usepackage{microtype}
\usepackage{fancyhdr}
\usepackage{enumitem}
\usepackage{hyperref}
\usepackage{lastpage}
\usepackage{parskip}

\hypersetup{colorlinks=true,urlcolor=black,linkcolor=black,
  pdftitle={STA/STB-603: Elements of Stochastic Processes -- BHU PYQ},pdfauthor={Banaras Hindu University}}

\pagestyle{fancy}\fancyhf{}
\renewcommand{\headrulewidth}{0.4pt}\renewcommand{\footrulewidth}{0.4pt}
\fancyhead[L]{\small   \textit{Banaras Hindu University}}
\fancyhead[R]{\small   \textit{STA/STB-603: Elements of Stochastic Processes}}
\fancyfoot[C]{\small Page  \thepage\ of \pageref{LastPage}}


\newlist{parts}{enumerate}{2}
\setlist[parts,1]{label=(\alph*),leftmargin=*,itemsep=5pt,topsep=3pt}
\setlist[parts,2]{label=(\roman*),leftmargin=*,itemsep=3pt,topsep=2pt}

\newcommand{\pts}[1]{\hfill   \textit{[#1]}}


\begin{document}


\begin{center}
  {\large\bfseries BANARAS HINDU UNIVERSITY}\\[2pt]
  {
\normalsize B.A./B.Sc. (Hons.) Semester VI Examination, 2022-23}\\[6pt]
  \rule{0.8\linewidth}{0.5pt}\\[6pt]
  {\large\bfseries Paper No. STA/STB-603: Elements of Stochastic Processes}\\[4pt]
  {
\normalsize Time: 3 Hours\hfill Maximum Marks: 70}
\end{center}
\rule{\linewidth}{0.4pt}
\smallskip



\noindent   \textit{Instructions: Answer any   \textbf{five} questions. All questions carry equal marks. Symbols have their usual meanings.}
\bigskip

\medskip


\noindent   \textbf{Question 1.}

\begin{parts}
  \item What do you mean by a random (stochastic) process? Also describe the classification of random processes according to state space and time domain. \pts{7}
  \item Explain the various classifications of states of a Markov chain (communicating, accessible, transient, recurrent, periodic, aperiodic, ergodic). \pts{7}
\end{parts}

\medskip\rule{\linewidth}{0.2pt}

\medskip


\noindent   \textbf{Question 2.}

\begin{parts}
  \item Define a wide-sense stationary (WSS) process with an example. If a process is strict-sense stationary (SSS), show that its mean is constant and its autocorrelation depends only on the time lag. \pts{7}
  \item Show that the process $X(t) = A\cos\lambda t + B\sin\lambda t$, where $A$ and $B$ are random variables, is wide-sense stationary if:
  
\begin{parts}
    \item $E(A) = E(B) = 0$
    \item $E(A^2) = E(B^2)$
    \item $E(AB) = 0$
  \end{parts}\pts{7}
\end{parts}

\medskip\rule{\linewidth}{0.2pt}

\medskip


\noindent   \textbf{Question 3.}

\begin{parts}
  \item Explain the terms state space and parameter space. Discuss with suitable examples where both state space and parameter space are discrete. \pts{7}
  \item Consider a time-homogeneous Markov chain with state space $S = \{0, 1, 2\}$ and transition matrix:
  \[
    P = 
\begin{pmatrix} \frac{1}{2} & \frac{1}{4} & \frac{1}{4} \\[4pt] \frac{1}{3} & \frac{1}{3} & \frac{1}{3} \\[4pt] \frac{1}{4} & \frac{1}{4} & \frac{1}{2} \end{pmatrix}
  \]
  
\begin{parts}
    \item Verify that each row sums to 1 and the matrix represents a valid transition matrix.
    \item Draw the transition graph for this process.
  \end{parts}\pts{7}
\end{parts}

\medskip\rule{\linewidth}{0.2pt}

\medskip


\noindent   \textbf{Question 4.}

\begin{parts}
  \item Let $X_n$ be a random variable representing the weather of a particular place on a given day: $X_n = 0$ if the day is rainy and $X_n = 1$ if the day is sunny. The one-step transition probability matrix is:
  \[
    P = 
\begin{pmatrix} 0.6 & 0.4 \\ 0.2 & 0.8 \end{pmatrix}
  \]
  Find the weather of the day in the distant future (steady-state distribution). \pts{7}
  \item Three boys A, B, and C are throwing a ball to each other. A always throws to B and B always throws to C, but C is equally likely to throw to B or to A. Show that the process is Markovian. Find the transition probability matrix and classify the states. \pts{7}
\end{parts}

\medskip\rule{\linewidth}{0.2pt}

\medskip


\noindent   \textbf{Question 5.}

Derive the formula associated with the Gambler's Ruin Problem. A gambler starts with initial capital of Rs.\ 12, and his probability of winning each round is $0.35$. His opponent starts with Rs.\ 8. Find:

\begin{parts}
  \item The probability of ruin for the first gambler.
  \item The expected duration of the game.
  \item Also find the expected duration of the game when the probability of winning is $0.5$ for each gambler.
\end{parts}\pts{14}

\medskip\rule{\linewidth}{0.2pt}

\medskip


\noindent   \textbf{Question 6.}

\begin{parts}
  \item Define a Poisson process. Show that the difference of two independent Poisson processes is not a Poisson process. \pts{7}
  \item Describe the time-dependent Poisson process. Also find its mean and standard deviation. \pts{7}
\end{parts}

\medskip\rule{\linewidth}{0.2pt}

\medskip


\noindent   \textbf{Question 7.}

Discuss the linear birth process (Yule--Furry process) and the pure birth process, with their assumptions and uses. \pts{14}

\medskip\rule{\linewidth}{0.2pt}

\medskip


\noindent   \textbf{Question 8.} Define the following terms:\pts{14}

\begin{parts}
  \item Aperiodic chain
  \item $n$-step transition probability
  \item Persistent (recurrent) state
  \item Periodicity of a state
  \item Ergodicity
  \item Probability of extinction in a branching process
  \item Cross-correlation function
\end{parts}

\vfill
\rule{\linewidth}{0.4pt}

\begin{center}\small
  \href{https://bhu-exam.vercel.app/}{Click here to visit our website}\\[4pt]
  \href{https://me-aryan.vercel.app/}{Click here to visit our contributor}\\[4pt]
  \href{https://quashan-me.vercel.app/}{Click here to visit our contributor}
\end{center}

\end{document}